% !TEX TS-program = pdflatex
% Lettre de motivation — Ingénieur(e) Sécurité, Centre des Opérations de Sécurité (SOC) — UIB
\documentclass[a4paper,11pt]{article}
\usepackage[utf8]{inputenc}
\usepackage[T1]{fontenc}
\usepackage[scaled=0.94]{helvet}
\renewcommand{\familydefault}{\sfdefault}
\usepackage[left=2.5cm,right=2.5cm,top=2cm,bottom=2cm]{geometry}
\usepackage{xcolor}
\usepackage[hidelinks]{hyperref}
\definecolor{navy}{RGB}{23,54,93}
\pagestyle{empty}
\setlength{\parindent}{0pt}
\setlength{\parskip}{0pt}
\hypersetup{pdftitle={Lettre de motivation - Baha Eddine Belhaj Mouldi - UIB Ingenieur Securite SOC},pdfauthor={Baha Eddine Belhaj Mouldi}}
\begin{document}

\begin{flushright}
{\bfseries Baha Eddine Belhaj Mouldi}\\
Tunis, Tunisie\\
+216 55 128 982\\
\href{mailto:bahaeddine.belhajmouldi@gmail.com}{bahaeddine.belhajmouldi@gmail.com}\\
\href{https://linkedin.com/in/bahamouldi}{linkedin.com/in/bahamouldi}\\
\href{https://bahamouldi.github.io/portfolio/}{Portfolio}
\end{flushright}

\vspace{0.8cm}

\begin{flushleft}
Tunis, le 21 août 2026

\vspace{0.8cm}

{\bfseries À l'attention de la Direction des Ressources Humaines --- UIB}

\vspace{0.6cm}

{\bfseries Objet : Candidature au poste d'Ingénieur(e) Sécurité --- Centre des Opérations de Sécurité (SOC)}
\end{flushleft}

\vspace{0.6cm}

Madame, Monsieur,

\vspace{0.4cm}

Diplômé Ingénieur en Génie Informatique de l'École Nationale d'Ingénieurs de Tunis (ENIT), je vous soumets ma candidature pour le poste d'Ingénieur(e) Sécurité au sein de votre Centre des Opérations de Sécurité. Le secteur bancaire, avec ses exigences élevées de disponibilité, de conformité et de confidentialité, correspond exactement au niveau de rigueur que j'ai recherché dans chacune de mes expériences en cybersécurité.

Au cours de mon projet de fin d'études chez Beehive Entreprises, j'ai conçu BeeWAF, une plateforme de sécurité applicative d'entreprise combinant 10 041 règles de détection et 3 modèles de machine learning, ainsi que des briques d'identité et d'accès (PKI, MFA, PAM) et 27 modules de sécurité (DLP, anti-DDoS, virtual patching). Ce projet a été audité et mis en conformité avec 7 référentiels reconnus (OWASP, PCI DSS, GDPR, NIST, ISO 27001, CIS, SOC 2), avec un mapping complet des techniques MITRE ATT\&CK et attribution d'acteurs de la menace.

En parallèle, dans le cadre de mes missions freelance chez Beehive Entreprises, j'ai administré et durci des environnements Windows Server / Active Directory et Linux, configuré et supervisé des VPN, proxy et règles de firewall réseau, et mené des tests d'intrusion sur 8 applications web et API avec rédaction de rapports de remédiation. Mon stage chez Streamlink m'a permis de construire un pipeline de détection et réponse automatisé pour la supervision SAP (Python/PyRFC, ELK, orchestration N8N de type SOAR, gestion des incidents sous TheHive), incluant la gestion des accès à privilèges et le contrôle de la Segregation of Duties. Mon stage chez Tunisie Telecom a complété cette expérience par un audit de posture réseau et un projet de détection orienté SOC (Elastic Stack, Kibana, règles Sigma).

Cette combinaison d'opérations SOC (triage d'alertes, corrélation de logs, investigation d'incidents), de sécurité systèmes et réseau (Windows Server, Active Directory, VPN, Proxy, Firewall) et de gestion des identités et accès (PKI, MFA, PAM, IAM) constitue selon moi une base solide pour contribuer efficacement à votre centre des opérations de sécurité, tout en continuant à approfondir mes compétences au contact de vos équipes.

Je me tiens à votre disposition pour un entretien et vous remercie de l'attention portée à ma candidature.

\vspace{0.8cm}

\begin{flushright}
{\bfseries Baha Eddine Belhaj Mouldi}
\end{flushright}

\end{document}
